\chapter{Literature Review}\label{chap:litreview}

\section{Clinical definitions and epidemiological context}
Sepsis-associated acute kidney injury (SA\mbox{-}AKI) has been characterized as the concurrence of sepsis and acute kidney injury with persistently poor outcomes. Sepsis has been defined by the Sepsis\mbox{-}3 consensus as life-threatening organ dysfunction arising from a dysregulated host response to infection, and the construct has been operationalized as suspected or documented infection together with an acute rise in organ failure, commonly a SOFA increase of at least two points \citep{Singer_2016}. Acute kidney injury has been defined by KDIGO through serum creatinine and urine output thresholds, in particular an absolute creatinine increase of at least $0.3\,\mathrm{mg/dL}$ within forty-eight hours, or at least $1.5\times$ baseline within seven days, or urine output of $0.5\,\mathrm{mL/kg/h}$ or less for six hours \citep{KDIGO_2012}. Many studies have labeled SA\mbox{-}AKI when AKI develops within a short interval after sepsis onset, frequently within seven days, to attribute renal dysfunction to the septic process rather than to competing causes \citep{Peerapornratana_2019,Poston_Koyner_2019}. The shift toward Sepsis\mbox{-}3 and KDIGO has reduced heterogeneity relative to older criteria and has improved comparability across cohorts \citep{Singer_2016,KDIGO_2012}.

The burden of SA\mbox{-}AKI in intensive care units has been described as substantial. It has been observed that AKI develops in roughly fifty to sixty percent of septic ICU patients, and that sepsis accounts for about forty to seventy percent of AKI in that setting \citep{Peerapornratana_2019,Poston_Koyner_2019}. When AKI complicates sepsis, marked increases in in-hospital mortality have been reported, with earlier analyses suggesting a six to eight fold rise compared with sepsis without AKI and mortality rates commonly between forty and sixty percent \citep{Poston_Koyner_2019,Peerapornratana_2019}. Prolonged ICU stays and an elevated risk of subsequent chronic kidney disease have also been recorded among survivors \citep{Poston_Koyner_2019,Peerapornratana_2019}. These epidemiological patterns have motivated the development of early risk stratification tools.

\section{Conventional severity indices and early warning systems}
General ICU indices such as SOFA, SAPS~II, and APACHE~II have been evaluated for SA\mbox{-}AKI outcomes and have tended to show only modest discrimination for short-term mortality or kidney-specific endpoints in this subgroup \citep{Luo_2022,Li_2023_SciRep,Qu_2024}. Classical regression frameworks, including logistic regression and Cox models, have been employed repeatedly, yet difficulties have been reported in capturing the nonlinear and interacting relationships among infectious physiology, hemodynamics, renal injury, and treatment exposures. As a result, performance has often not transported reliably across hospitals \citep{Kang_Yoon_2023,Luo_2022,Li_2023_SciRep}. Earlier variability in sepsis and AKI definitions added to the challenge by complicating cohort selection and cross-study comparison prior to widespread adoption of Sepsis\mbox{-}3 and KDIGO \citep{Singer_2016,KDIGO_2012}.

\section{Datasets and cohort construction in SA-AKI modeling}
Large electronic health record repositories have been used to develop prognostic models for SA\mbox{-}AKI, and MIMIC\mbox{-}IV has been the most common source for cohort construction and feature derivation \citep{Johnson_2023_MIMIC}. Studies have typically restricted analysis to the first ICU admission per patient, excluded patients younger than eighteen years, and removed ICU stays shorter than forty-eight hours. A central design decision has involved the temporal window for feature construction. Fixed early windows, for example the first twenty-four or forty-eight hours after admission or sepsis recognition, have been used to support timeliness, whereas whole-stay aggregation from admission to discharge has also appeared frequently and has created a risk that future information relative to an intended decision point was introduced into model inputs.

\section{Feature engineering, selection, and leakage risks}
The handling of high-dimensional clinical variables has relied on screening by correlation or variance inflation factor, wrapper methods such as recursive feature elimination, and embedded strategies that exploit model-based importance estimates. It has been common to describe feature selection steps that were informed by the full dataset. When selection has been conducted in this manner, the final predictor set has been influenced by outcomes from records that were later evaluated, and optimism in reported discrimination has been encouraged \citep{Kang_Yoon_2023}. Claims that VIF-based pruning enhances stability and generalizability have also appeared without pre–post diagnostics or external tests, and the relevance of VIF is limited for classes of models that tolerate correlated predictors. In several studies, length of stay has been included among features for mortality prediction. Because time to event is strongly associated with whether death has occurred during the admission, that inclusion has acted as a conduit for post-baseline information and has dominated other covariates during training. Redundancy among variables has also been noted. For example, total urine output and average urine output become scalar multiples under a fixed observation horizon, and diverging importance profiles for those variables have suggested that non fixed or whole-stay windows were used, which reduces alignment with early prediction. Dependencies between creatinine and derived eGFR have similarly complicated interpretation when both were included without a clear rationale.

\section{Missing data mechanisms and imputation choices}
The mechanisms that govern missingness have rarely been examined in depth in this literature. Assumptions that data were missing completely at random or missing at random have been stated without empirical assessment, and simple percentage-based rules have been used to trigger imputation. Mean imputation for continuous variables has been a common choice even for biomarkers that show marked skew, such as serum creatinine, which has been described as right skewed and well approximated by a log-normal distribution \citep{Haeckel_2010}. Multiple imputation by chained equations and random forest variants has been widely used, and in many cases those imputation models have been fitted across the entire dataset prior to train–test separation. Because chained equations use all observed variables to impute each target variable, global fitting has allowed information from records later treated as held out to influence imputed values, and optimistic estimates of model performance have been encouraged.

\section{Class imbalance, evaluation strategy, and thresholding}
Mortality among SA\mbox{-}AKI cohorts has been imbalanced, and class imbalance has been handled most often by oversampling techniques, including SMOTE and ADASYN \citep{Chawla_2002,He_2008}. The two methods have occasionally been applied together even though ADASYN is itself a variant that emphasizes synthesis in hard regions of the feature space. It has been reported that SMOTE may generate synthetic points in overlapped or noisy areas and that performance can degrade in high-dimensional settings unless additional filtering is applied \citep{Blagus_2013,Saez_2015}. Concentration of synthesis near low-density or borderline regions by ADASYN can shift decision boundaries around outliers and may harm robustness under severe imbalance or distribution shift \citep{Krawczyk_2016}. Comparisons with alternative imbalance strategies, for instance through intrinsic class weighting or hybrid sampling, have not always been provided. Evaluation strategies have often relied on a single split of the data into training and test sets without a separate validation split, and explicit calibration analyses and principled threshold selection have been omitted in several reports.

\section{Model classes, tuning, calibration, and interpretability}
Tree-based ensembles, notably extreme gradient boosting, have been adopted widely in this domain to represent nonlinearities and interactions among clinical variables, and stronger internal discrimination than baseline ICU scores has often been reported \citep{Chen_Guestrin_2016,Luo_2022,Li_2023_SciRep,Qu_2024,Wang_2024,Chen_2025,Le_2025}. Hyperparameter tuning has usually been performed through grid search cross validation, and claims of optimality have been made despite the presence of complex, nonconvex search spaces in which only local minima are guaranteed by such procedures. Calibration has been reported inconsistently, and explanations based on Shapley values have been used to describe feature contributions, although stability of such explanations has not always been explored. In several comparisons, confidence intervals for ensembles and penalized linear models have overlapped, and differences attributed to model class have been entangled with choices about temporal windows, selection scope, and imputation design.

\section{External validity, transportability, and deployment timing}
Generalization beyond the MIMIC\mbox{-}IV development site has remained a consistent concern. MIMIC\mbox{-}IV reflects a single tertiary-care center in Boston with local care patterns and population structure that are not broadly representative \citep{Johnson_2023_MIMIC}. External evaluations of ICU prediction models have reported loss of accuracy when transported across hospitals, suggesting the presence of center-specific differences that matter for performance \citep{Rockenschaub_2024}. Assertions about applicability to diverse populations have therefore required cautious interpretation when external testing has not been provided. In addition, the timing of features relative to the intended decision point has often been unspecified, and whole-stay aggregation has created ambiguity about whether a model could power actionable alerts at twenty-four or forty-eight hours.

\section{Detailed appraisal of representative studies}

\subsection{Chen et al.\ (2025): SA-AKI mortality in MIMIC\mbox{-}IV}
A recent contribution by Chen et al.\ reported an AUROC of 0.878 for mortality prediction in a SA\mbox{-}AKI cohort drawn from MIMIC\mbox{-}IV, and the result was presented as an improvement over the AUROC 0.794 attributed to Li et al.\ \citep{Chen_2025,Li_2023_SciRep}. The cohort retained only the first ICU admission per patient and excluded patients younger than eighteen years and ICU stays shorter than forty-eight hours. The modeling pipeline combined VIF-based filtering with recursive feature elimination to produce a set of twenty-four predictors, after which several variables were reintroduced as expert inputs without quantitative justification. Because selection was described as having been performed on the entire dataset, the final predictor set was influenced by outcomes from records later treated as held out and the risk of information leakage was created \citep{Kang_Yoon_2023}. A claim of generalizability across diverse populations was made, yet MIMIC\mbox{-}IV represents a single center and patient mix, and performance drift across hospitals has been documented \citep{Johnson_2023_MIMIC,Rockenschaub_2024}. The introduction criticized Li et al.\ for relying on forty-four variables on the grounds that the set might be too narrow to capture the complexity of SA\mbox{-}AKI, although the final model by Chen et al.\ used an even smaller set of twenty-four features, which created an internal inconsistency in the narrative \citep{Chen_2025,Li_2023_SciRep}. Hyperparameters were tuned with grid search cross validation and described as optimal, although such procedures identify settings within a discrete grid and do not ensure global optimality in nonconvex spaces. Claims of timely, accurate, and actionable interventions were advanced, yet the temporal window for aggregating time-varying biomarkers was not specified. Because Sepsis\mbox{-}3 and KDIGO anchor observation to early horizons, an early extraction window would have been expected for actionable deployment; multiple signals suggested that whole-stay aggregation may have been used, which would introduce future information relative to early prediction. Missing values were imputed by means that were tied to percentage thresholds, and mean imputation was applied for variables with up to twenty percent missingness. This approach is not well aligned with skewed distributions such as serum creatinine, which has been described as right skewed and is well modeled by a log-normal distribution \citep{Haeckel_2010}. Missingness indicators were not created and the mechanism of missing data was assumed rather than studied. VIF filtering was described as enhancing stability and generalizability without pre–post diagnostics or external assessment and has limited relevance for ensembles that tolerate correlated predictors. The base estimator used within recursive feature elimination was not stated, and a linear estimator would impose a separability assumption during selection. Length of stay was included among predictors, and this inclusion introduced a strong proxy for time to event that is not available prospectively for early prediction. The simultaneous inclusion of total and average urine output suggested redundancy if a fixed window had been used, and differences in importance between the two variables implied that a non fixed or whole-stay window was applied. The evaluation strategy used a seventy-five to twenty-five train–test split without a separate validation split, and explicit calibration and threshold selection were not described. Class imbalance was handled by applying SMOTE and ADASYN together, a combination that lacks a clear rationale and that may synthesize examples in overlapped or noisy regions or near outliers and thereby distort decision boundaries \citep{Chawla_2002,He_2008,Blagus_2013,Saez_2015,Krawczyk_2016}. Univariate comparisons between survivors and a group termed non-survivors were presented with t tests even though laboratory variables such as creatinine generally deviate from normality \citep{Haeckel_2010}. The reported confidence intervals for ensemble and linear models overlapped, with AUROC 0.878 (95\% CI 0.859 to 0.897) for extreme gradient boosting and AUROC 0.849 (95\% CI 0.824 to 0.873) for logistic regression, which supported the interpretation that any advantage might be smaller than suggested once temporal alignment and leakage are addressed \citep{Chen_2025}.

\subsection{Li et al.\ (2023): SA-AKI mortality in MIMIC\mbox{-}IV}
Li et al.\ assembled a SA\mbox{-}AKI cohort from MIMIC\mbox{-}IV by retaining only the first ICU admission and excluding ICU stays shorter than forty-eight hours. Features were extracted from the first twenty-four hours and were aggregated using maximum values, which provided a clear early window for model inputs \citep{Li_2023_SciRep}. Lasso-regularized logistic regression was used for feature selection, and selection was described across the full dataset rather than within development folds. Because selection then involved outcomes from records later used for evaluation, optimism in reported performance was encouraged. The linear framework also prioritized main effects and could miss interactions and subtle nonlinear patterns that tree ensembles tend to capture. Multiple imputation by chained equations with random forest implementations was applied globally across the dataset rather than learned on development data and then applied to evaluation data, and missingness indicators were not constructed. The data were split into training and validation in an eighty to twenty ratio, and a discrimination gap favoring extreme gradient boosting over logistic regression was reported with AUROC 0.794 versus 0.730, although the magnitude of this gap remained difficult to interpret in light of the selection and imputation scope \citep{Li_2023_SciRep}.

\subsection{Hu et al.\ (2021): survival analysis in MIMIC\mbox{-}III}
A retrospective analysis by Hu et al.\ used MIMIC\mbox{-}III with smaller sample size and older coding structure to study SA\mbox{-}AKI outcomes \citep{Hu_2021}. The cohort retained only the first ICU admission and excluded patients younger than eighteen years and ICU stays shorter than forty-eight hours, and rows with more than five percent missing variables were removed. Lasso regression and Cox regression were used, and a C index of 0.72 was reported on a validation set. Biomarkers were aggregated by medians across the entire admission, therefore features included information that would not be available at an early decision point. Features with more than twenty percent missingness were dropped without investigation of mechanisms, and multiple imputation by chained equations was performed on the full dataset prior to splitting, which created leakage. The test set contained one hundred and two rows, and reliability of performance estimates was limited by that size. Ethnicity was used as a predictor, which raised concerns about representativeness and fairness for transportability. Survival models at seven, fourteen, and twenty-eight days yielded AUCs of 0.77, 0.73, and 0.71 respectively.

\subsection{Roknaldin et al.\ (2024): sepsis to AKI prediction in MIMIC\mbox{-}III}
A recent preprint constructed models that predict AKI among septic ICU patients using MIMIC\mbox{-}III \citep{Roknaldin_2024_arXiv}. Adults aged eighteen to eighty-nine years with sepsis by ICD\mbox{-}9 codes were included after excluding multiple admissions, ICU stays shorter than forty-eight hours, and records with more than twenty percent missingness. Features from the first twenty-four hours covered demographics, comorbidities, vitals, laboratory results, and early interventions. Multiple imputation was used for missing values, a correlation filter retained twenty-three predictors, and class imbalance was addressed with SMOTE. Six algorithms were compared, and logistic regression achieved the strongest internal results with test AUC 0.887 and validation AUC 0.857 together with a Brier score near 0.13 and acceptable isotonic calibration. Top predictors included urine output, bilirubin measures, blood urea nitrogen, minimum eGFR, weight, and indicators of vasopressor use and mechanical ventilation. The transparency of the feature list and the early window supported interpretability, while several design choices encouraged optimism. Feature selection and SMOTE were described prior to the final split into development and evaluation, so that the set of predictors and synthetic examples were influenced by records that later appeared in evaluation. Multiple imputation also appeared to have been fitted globally. A fixed correlation threshold risked removal of variables that are weakly correlated but jointly informative, and the inclusion of minimum and maximum eGFR alongside creatinine introduced redundancy because eGFR is derived from creatinine and demographics. For descriptive group comparisons, t tests were used for markers that are typically skewed, and the allocation of half the data to validation and one tenth to testing left a small test set. Calibration procedures were not tied to a recalibration step that excluded the test cohort, which limited the strength of calibration claims on that set.

\subsection{Chaudhary et al.\ (2020): phenotype discovery in MIMIC\mbox{-}III}
An unsupervised study in \emph{CJASN} sought to identify SA\mbox{-}AKI phenotypes rather than predict mortality directly \citep{Chaudhary_2020_CJASN}. An autoencoder was used to learn representations that were then clustered into three subgroups with distinct twenty-eight day mortality and dialysis rates. Differences were driven by bilirubin, lactate, transaminases, and inflammatory counts, and the highest risk subgroup showed markedly higher dialysis use and death. The analysis relied on whole-stay trajectories of laboratory and vital signs, which embeds information not available at early decision times. The sepsis cohort was constructed by codes, which can misalign the timing of the septic insult, and external validation of cluster stability was not performed. The exploratory links between clusters and treatments were not framed as causal and would require dedicated designs to support treatment guidance.

\subsection{Lin et al.\ (2025): progression to chronic kidney disease in MIMIC\mbox{-}IV}
A study in \emph{Informatics in Medicine Unlocked} examined progression from SA\mbox{-}AKI to incident chronic kidney disease using MIMIC\mbox{-}IV \citep{Lin_2025_IMU}. Extreme gradient boosting with SHAP explanations was trained on features collected during the first twenty-four hours, and high internal discrimination was reported. Thresholds for age, maximum creatinine, minimum hemoglobin, and lactate were described, and a feature stability analysis across AKI stages was presented. Although the early window and the inclusion of decision-curve analysis aided interpretability, variable screening by Lasso and stepwise logistic regression prior to final modeling, if performed across the full dataset, would have allowed outcomes from evaluation records to influence the screened set. Outcome assessment relied on follow up up to five years, and selection by loss to follow up as well as center-specific practice patterns could have influenced estimates. Thresholds based on SHAP dependence plots have known sensitivity to sampling variation, and replication would be expected for robust adoption.

\subsection{He et al.\ (2021): acute kidney disease after SA\mbox{-}AKI}
An article in \emph{Frontiers in Medicine} predicted subsequent acute kidney disease rather than in-hospital mortality \citep{He_2021_FrontMed}. A single-center development cohort was used, and external validation was undertaken in MIMIC\mbox{-}III. Recurrent neural networks with long short-term memory units were compared with a decision tree and logistic regression. The recurrent model achieved an internal AUROC of 1.00 and strong performance in the external cohort, and the change in nonrenal SOFA between day one and day three emerged as an important signal. The development cohort contained two hundred and nine patients, which is a modest sample for a deep architecture with more than two dozen features, and perfect discrimination in such data is difficult to reconcile with typical ICU noise and suggests that overfitting, simplified labels, or other forms of information leakage may have been present. The dynamic window implied by the day one to day three delta improves clinical timing, while the time of alert generation relative to clinical workflow was not specified. It was not always clear that imputation and preprocessing steps were confined to development data.

\subsection{Kantola et al.\ (2025): admission procalcitonin and risk of AKI}
A multicenter retrospective cohort across three hospitals evaluated admission procalcitonin for prediction of subsequent AKI and reported an odds ratio of approximately 1.01 per ng/mL with poor overall discrimination near ROC AUC 0.59 \citep{Kantola_2025}. The heterogeneity of practical cutoffs across prior work and the inconsistent relationships between procalcitonin levels and AKI severity were emphasized, and the conclusion favored limited standalone clinical utility for early risk stratification. Procalcitonin is not available in MIMIC\mbox{-}IV, which restricts replication in that database.

\section{Synthesis and implications for this thesis}
The literature reviewed above sketches a consistent picture in which window definition, the scope of data-dependent preprocessing, handling of missingness and imbalance, and the availability of external evidence play decisive roles in the credibility of reported scores. Studies that relied on whole-stay aggregation, on inclusion of length of stay or other post-baseline correlates, and on global feature selection or imputation have tended to report stronger discrimination, yet those gains have been paired with ambiguity about practical timing and with higher risk of optimism. When early windows were enforced clearly, performance remained competitive and the path to bedside use became clearer. Several reports presented overlapping confidence intervals between ensembles and penalized linear models, which suggests that differences in pipeline design have often overshadowed differences in the learning algorithms themselves \citep{Kang_Yoon_2023}. Generalization beyond a single tertiary-care center has remained a key uncertainty, and drift across hospitals has been documented \citep{Johnson_2023_MIMIC,Rockenschaub_2024}. In this thesis, these observations are treated as design constraints for methodological choices, and the subsequent chapters are organized to respond to the gaps that were identified here while keeping the emphasis on temporally anchored modeling, leakage-aware preprocessing, calibrated evaluation, and transparency about what information is and is not available at the moment when a clinical decision is intended to be supported.
